\chapter{Positive Current}

\begin{thebibliography}{n}
\bibitem[Rud]{Rud}
W. Rudin. \textit{Functional analysis.} 2nd edn. International Series in Pure and Applied Mathematics. McGraw-Hill, Inc., New York. (1991.)
\bibitem[Rud 2]{Rud2}
W. Rudin. \textit{Real and complex analysis. } 3rd ed. New York, NY: McGraw-Hill 
\bibitem[NO]{NO}
J. Noguchi, T.Ochiai \textit{Geometric Function Theory in Several Complex Variables} Translations of Mathematical Monographs
Volume: 80; 1990; 282 pp
\end{thebibliography}
 
\cite[Chapter 3]{NO}を参考にしている. 

\section{}

\begin{tcolorbox}[mybox]
\begin{prop}
 \(M\) 向きづけ可能 \(C^\infty\)多様体. 
  \(\dim_{\mathbb R} M = m\),  \(0 \le p \le m\)とする. 
\[
[\,\cdot\,] : \mathcal{L}^{m-p}_{\mathrm{loc}}(M) \to \mathcal{K}'_p(M)\ (\subset \mathcal{D}'_p(M))
\quad
\omega \longmapsto \left(\varphi \mapsto \int_M \omega \wedge \varphi\right)
\]
は全射となる. 
\end{prop}
\end{tcolorbox}

\begin{rem}
 \(\exists f:\mathbb{R}\to\mathbb{R}\) injective \(C^\infty\)関数と \(\exists A \subset \mathbb{R}\) Lebesgue可測なもので. 
\(f^{-1}(A)\) がLebesgue可測にならないものがある.

\(\exists f:[0,1]\to[0,1]\) 同相写像, \(\exists A \subset \mathbb{R}\) Lebesgue可測なもので. 
\(f^{-1}(A)\) がLebesgue可測にならないものがある.
\end{rem}
\begin{ques}
正則写像やalgebraic mapの場合はどうなるか？
\end{ques}

\begin{proof}
% [page 1]
 \([\,\cdot\,]: \mathcal{L}^{m-p}_{\mathrm{loc}}(\bullet) \to \mathcal{K}'_p(\bullet)\)がsheafの射なので,  \(M = U \subset \mathbb{R}^m\)という開集合にして良い. 
さらに \(\omega = \sum_J \omega_J dx^J \in L^{m-k}_{\mathrm{loc}}(U)\)とする．
\footnote{ただし$ L^{m-k}$はalmost everywhere で同一視する前のもので, 係数は関数とする. (つまり同一視する前のもの)}

任意の$\varphi \in \mathcal{K}^k(U)$について, \(\int_U \omega \wedge \varphi = 0\)であると仮定する. 
示すことは任意の添字$J$について,  \(\omega_J = 0\) である. 

任意の \(J=(i_1,\cdots,i_{m-k})\)について, 
 \(J^\circ := (j_1,\cdots,j_k)\) を
\[
\{i_1,\cdots,i_{m-k},\, j_1,\cdots,j_k\}=\{1,\cdots,m\}.
\]
となるようにとる. 任意の\(\varphi \in \mathcal{K}(U)\)について, 
\[
0=
\int_U \omega \wedge \underbrace{\varphi\, dx^{J^\circ}}_{\in \mathcal{K}^k(U)}
   = \int_U \omega_J \varphi\, dx^J \wedge dx^{J^\circ}.
\]
よって, 任意の\(\forall \varphi \in \mathcal{K}(U)\)について, 
\(\int_U \omega_J \varphi\, dm = 0\)である. ここで$dm$をLesbegue 測度とする. 
よって次の補題からいえた. 
\end{proof}

\begin{tcolorbox}[mybox]
\begin{lem}
\(U \subset \mathbb{R}^m\) 開集合, \(f \in L_{\mathrm{loc}}(U)\)とする. 
今任意の$\varphi \in \mathcal{K}(U)$について. 
\(
\int_U f\varphi\, dm = 0 
\)
が成り立つならば,  測度\(m\)についてalmost everywhereで\(f=0\)である.
\end{lem}
\end{tcolorbox}
$ \mathcal{K}(U)$
は$U$上のコンパクトサポートを持つ連続$\C$値関数の集合である.
\begin{proof}

[Step 1.] 
\(f \in L^1(U)\) (つまり\(\int_U |f|\,dm < \infty\))であることを仮定して良い. 

なぜならば \(U=\bigcup_{i\geq 1} U_i\) 可算個の開被覆で
\(\overline{U_i}\) がコンパクトとなるものをとる. 
 \(\{\varphi_i\}_{i\geq 1}\) を１の分割とする. (Supp \(\varphi_i \subset U_i\)となる)

すると\(\varphi_i f \in L^1(U)\) かつ \(\forall \varphi \in \mathcal{K}(U)\)について
\[
\int_U (\varphi_i f)\varphi\,dm
=
\int_U f \underbrace{(\varphi_i\cdot \varphi)}_{\in \mathcal{K}(U)}\,dm
=
0
\]
よって, almost everywhereで$\varphi_i f=0$となるので, 
\[
f=\sum_{i\geq 1}\varphi_i f = 0 \ \text{a.e.}
\]

[Step 2] 
\(\quad f \in L^1(U)\) の場合に示す. 

%Idea.\[\mathcal{K}(U) \to \mathbb{C};\ \varphi \mapsto \int_U f\varphi\,dm\]
%\textcolor{red}{[囲み内注記]}
%これは supp narrow にいて bdd \textcolor{red}{\illegible{英字略記}}. Riesz rep
%\textcolor{red}{\illegible{日本語文}}

任意のLesbegue可測集合 \(E \subset U\) について, 
\[
\lambda(E)=\int_E f\,dm.
\]
とする. すると\(\lambda\)はcomplex measureになる.\footnote{$\lambda$がcomplex measureとは
$\lambda:\ \{\text{Leb meas set}\}\to \mathbb{C}$であって, 完全加法性
$\lambda\Bigl(\bigsqcup_{i=1}^{\infty} A_i\Bigr)
=
\sum_{i=1}^{\infty}\lambda(A_i)$が成り立つこと. ただし右辺の収束は絶対収束を要求する}

さて$\lambda$の total variationを
Lesbegue可測集合\(E\subset U\)について
\[
|\lambda|(E):=\sup\left\{\sum_{i=1}^{\infty}|\lambda(E_i)|\ \middle|\ E=\bigsqcup_{i=1}^{\infty}E_i,\ E_i\subset U \text{Lesbegue可測}\right\}
\]
とする. 
\cite[Chapter 6]{Rud2}より,  \(|\lambda|\) は finite positive measure, つまり$|\lambda|(U)<\infty$となる測度である. (この測度$|\lambda|$はtotal variation measureと呼ばれる. )

\cite[Theorem 6.12]{Rud2}よりRadon-Nikodymの定理から
ある\( h \in L^1(U,|\lambda|)\)であって, 
任意の$x \in U$について\(
|h(x)|=1\)かつ
任意のLesbegue可測集合$E\subset U$について
\[
\lambda(E)=\int_E h\,d|\lambda|.
\]
となるものが存在する. \footnote{このことを$d\lambda = h\,d|\lambda|$
と表し. \(\lambda\)のpolar expressionともいう}

このとき任意の\(\varphi \in \mathcal{K}(U)\)について. 
 \[\displaystyle \underbrace{\int_U \varphi h\,d|\lambda|}_{=:\int_U \varphi\,d\lambda} = \int_U \varphi f\,dm \underset{\text{仮定}}{=} 0\]
これを示す.

まず$E$をLesbegue可測集合で$\varphi=\chi_E$, つまり
$$\varphi=\chi_E,\qquad
\chi_E(x)=
\begin{cases}
1^\circ & x\in E\\
0 & x\notin E
\end{cases}
$$
である場合は
\[
\int_U \chi_E h\,d|\lambda|
\underset{\text{$\chi_E$の定義}}{=}
\int_E h\,d|\lambda|
\underset{\text{定義}}{=}
\lambda(E)
\underset{\text{定義}}{=}
\int_E f\,dm
\underset{\text{$\chi_E$の定義}}{=}
\int_U \chi_E f\,dm.
\]
であるのでOK.
\( \varphi\)が単関数(support関数$\chi_E$の有限な線型結合となる関数)の時もOK.

$ \varphi$が0以上の有界関数の時, 
単関数近似
\(s_i \to \varphi \)をとる. (ただし$s_i \le s_{i+1} \le \varphi $となるようにとる)
すると
\[
\int_U s_i h\,d|\lambda| \xrightarrow[i\to\infty]{} \int_U \varphi h\,d|\lambda|
\quad 
\int_U s_i f\,dm \xrightarrow[i\to\infty]{} \int_U \varphi f\,dm
\]
である, そして
\[
|s_i h|\leq \underbrace{|\varphi|}_{\text{bounded}}|h|,\qquad h\in L^1 \text{ w.r.t. } |\lambda|
\quad \text{and} \quad
|s_i f|\leq |\varphi||f| \in L^1 \text{ w.r.t. } m
\]
であるので, Lesbegue収束定理より$\int_U \varphi h\,d|\lambda|= \int_U \varphi f\,dm$となる. 
以上より$\varphi \in \mathcal{K}(U)$の時でもOKである. 

さて\(|\lambda|(U)=0\)ならば \(|\lambda|=0\) であることを示す. 
もしこれがいえたら\(\lambda=0\)が言えて,  任意のLesbegue可測集合 \( E\subset U\) について
\[
0\underset{\lambda=0}{=}\lambda(E)
\underset{\text{定義}}{=}\int_E f\,dm.
\]
となり, $f=0$ a.e.が言える.
今
\begin{align*}
|\lambda|(U)
&=\int_U d|\lambda|
\underset{|h|=1}{=}
\int_U \overline{h}\,h\,d|\lambda| \\
&\underset{\forall \varphi\in \mathcal{K}(U)}{=}\int_U (\overline{h}-\varphi)h\,d|\lambda|
\underset{}{=}\int_U \operatorname{Re}\bigl((\overline{h}-\varphi)h\bigr)\,d|\lambda| \\
&\underset{}{\le}
\int_U |(\overline{h}-\varphi)h|\,d|\lambda|\\
&\underset{|h|=1}{=}\int_U |\overline{h}-\varphi|\,d|\lambda|
\end{align*}
ここで$\int_U |\overline{h}-\varphi|\,d|\lambda|$は
\(\overline{h}\) と \(\varphi\) の \(L^1(U,|\lambda|)\)-距離である.

よって下のclaimより, $\overline{h} \in L^1(U,|\lambda|)$について
$\varphi_i \to \overline{h}$となる$\varphi_i \in \mathcal{K}(U)$をとれば, 
$|\lambda|(U)=0$がいえる. 

\begin{claim}
\(\mathcal{K}(U)\subset L^1(U,|\lambda|)\) は
 \(L^1\)-normでdenseである. 
\end{claim}
次の定理を思い出す. 
\begin{tcolorbox}[mybox]
\begin{thm}\cite[Theorem 3.14]{Rud2}
\(X\) locally compact Hausdorff,  \(\mathcal{M}\) $X$の\(\sigma\)-algebra, 
\(\mu\)を\(\mathcal{M}\)上のpositive measureとする. 

 \((X,\mathcal{M},\mu)\) が次を満たすとする.
\begin{itemize}
\item[(0)] \(\mathcal{M}\) は Borel algebraを含む. (つまり$\mathcal{M}$が$X$の全ての開集合を含む)
\item[(1)] 任意のコンパクト集合 \( K\subset X\) について \(\mu(K)<\infty\)
\item[(2)] (外部正則性) 任意の\( E\in\mathcal{M}\)について$\mu(E)=\inf\{\mu(V)\mid E\subset V\underset{\text{open}}{\subset} X \}\)
\item[(3)] (内部正則性) 任意の \(E\in\mathcal{M}\), \(E\) が開集合または\(\mu(E)<\infty\)ならば
\(
\mu(E)=\sup\{\mu(K)\mid K\underset{\text{cpt}}{\subset} E\}
\)
\item[(4)](完備性) 任意の\(E\in\mathcal{M},\ A\subset E,\ \mu(E)=0\)について, $A\in\mathcal{M}$
\end{itemize}
% [page 1]

この時 \(\{f:X\to\mathbb{C}\ \text{conti} \mid \operatorname{supp} f \text{ cpt}\}\subset L^1(X,\mu)\)
は\(L^1\)-normで
denseである. 
実は$L^p,\ 1\leq p<\infty$でも成り立つ. 
\end{thm}
\end{tcolorbox}

\(|\lambda|\) が定理の仮定を満たすことをみる. 
f\(|\lambda|\)は有限測度で, 
$ (0)\sim (3)$は OK(多少議論は必要, 有限Borel測度の正則性)
(4) については \((U,\text{Lesbegue可測集合},|\lambda|)\) の
測度空間としての完備かをとれば, 
\((0)\sim(3)\) を満たしながら, (4)も成立する.
この時$L^1$空間が大きくなるだけである. よってclaimが成り立つ. 
\end{proof}


\section{Positive currents}

\subsection{(a) Type of currents }

$\mathbb{C}$とかいたら pair  $(\mathbb{C}, i)$のこととする.
ここで$i$は$i^2=-1$となるfixした元とする. 

$m \in \mathbb{Z}_{\geq 1}$について, 
\[
\mathbb{C} = \mathbb{R}^2
\quad \ z=x+iy \leftrightarrow (x,y)
\]

\[
\mathbb{C}^m=\mathbb{R}^{2m}\quad 
(z^1,\dots,z^m) \leftrightarrow (x^1,y^1,x^2,y^2,\dots,x^m,y^m)
\]
と同一視する. 
$z^j=x^j+iy^j$とする. 

またcomplex manifoldは常に 2nd countable とする.

\begin{tcolorbox}[mybox]
\begin{defn}
\label{defn-M-4.1}
$ U \subset \mathbb{C}^m$をopenとするときに

\[
dz^j = dx^j + i\,dy^j,\quad d\bar{z}^j = dx^j - i\,dy^j
\]

\[
\frac{\partial}{\partial z^j} := \frac{1}{2}\left(\frac{\partial}{\partial x^j} - i\frac{\partial}{\partial y^j}\right),\quad
\frac{\partial}{\partial \bar{z}^j} := \frac{1}{2}\left(\frac{\partial}{\partial x^j} + i\frac{\partial}{\partial y^j}\right),
\]
\end{defn}
\end{tcolorbox}
次のnotationをつかう.

\begin{tcolorbox}[mybox]
\(k \geq 0\)について, 
\[
\{m; k \} := 
\left\{(j_1,\ldots,j_k)\in \mathbb{Z}^k \,\middle|\, 1 \leq j_1 < \cdots < j_k \leq m\right\}
\]

　任意の\(J=(j_1,\dots,j_k)\in \{m;k\}\)について, 
\[
dz^J := dz^{j_1}\wedge \cdots \wedge dz^{j_k}
\quad
d\bar{z}^{\bar{J}}:= d\bar{z}^{j_1}\wedge \cdots \wedge d\bar{z}^{j_k}.
\]
とする.  \(U\subset \mathbb{C}^m\) open, 
\(
 \varphi : U \to \Lambda^{\ell}T^{*}U\otimes_{\mathbb{R}} \mathbb{C}
 \)
となる
sectionについて, 
\[
\varphi = 
\sum_{\substack{p,q\geq 0\\ p+q=\ell}}
\sum_{J\in \{m;p\} , \bar{K}\in\{m;q\}}
\varphi_{J\bar{K}}\, dz^J\wedge d\bar{z}^{\bar{K}}
\]
となる唯一のmap\(\varphi_{J\bar{K}} : U \to \C\)が存在する. 
\end{tcolorbox}

\begin{tcolorbox}[mybox]
\begin{defn}
\(U \subset \mathbb{C}^m\) openについて次にように定める.
\begin{itemize}
\item 
$\mathcal{C}^{(p,q)}(U)
:=
\left\{
\sum_{J\in \{m;p\},  K\in \{m;q\}}
\varphi_{J\bar{K}}\, dz^J\wedge d\bar{z}^{\bar{K}}
\ \middle|\
\varphi_{J\bar{K}} : U \to \mathbb{C}\ \text{連続}
\right\}$
% [page 1]
\item $\mathcal{K}^{(p,q)}(U)=\mathcal{K}^{p+q}(U)\cap \mathcal{C}^{(p,q)}(U)$
\item $\mathcal{E}^{(p,q)}(U)=\mathcal{E}^{p+q}(U)\cap \mathcal{C}^{(p,q)}(U)$
\item $\mathcal{D}^{(p,q)}(U)=\mathcal{D}^{p+q}(U)\cap \mathcal{C}^{(p,q)}(U)$
\item 
$L^{(p,q)}_{\mathrm{loc}}(U):=
\left\{
\sum_{\substack{|J|=p\\|K|=q}}'
\varphi_{J\bar{K}}\,dz^J\wedge d\bar{z}^{\bar{K}}
\ \middle|\
\varphi_{J\bar{K}}\in L^1_{\mathrm{loc}}(U)
\right\}
$
\item 
$\mathcal{L}^{(p,q)}_{\mathrm{loc}}(U)
=
\mathrm{Im}\bigl(L^{(p,q)}_{\mathrm{loc}}(U)\to \mathcal{L}^{\,p+q}_{\mathrm{loc}}(U)\bigr)
$
とする. ただし$\mathcal{L}$はalmost everywhereの同値類で割った空間とする. 
\end{itemize}
\end{defn}
\end{tcolorbox}

\begin{rem}
\label{rem-M-4.3}
\[
\mathcal{K}^{k}(U)=\bigoplus_{p+q=k}\mathcal{K}^{(p,q)}(U),
\]
が成り立つ$\mathcal{E}^k,\ \mathcal{D}^k,\ \ldots\ $も同様である. 
\end{rem}

% [page 1]
\begin{tcolorbox}[mybox]
\begin{lem}
\label{lem-M-4.4}
\(\mathbb{C}^m \underset{\text{open}}{\supset} U \xrightarrow{\,f\,} V 
 \underset{\text{open}}{\subset} \mathbb{C}^n\)を正則写像とするとき, 
 \( J\in \{n;p\},\ \forall K\in \{n;q\},\)
について, 
\(
f^{*}(dz^{J}\wedge d\bar{z}^{K}) 
\)も type (p,q)である. 
\end{lem}
\end{tcolorbox}

\begin{proof}
\begin{align*}
f^{*}dz^{j}=f^{*}(dx^{j}+i\,dy^{j}) \underset{Cauchy-Riemann}{=}\
=\sum_{k=1}^{m}\frac{\partial f^{j}}{\partial w^{k}}\,dw^{k}
\end{align*}
同様にして, 
$
f^{*}d\bar{z}^{j}=
\sum_{k=1}^{m}\frac{\partial \overline{f^{j}}}{\partial \overline{w}^{\,k}}\,d\overline{w}^{\,k}$となるのでいえる. 
\end{proof}

\begin{tcolorbox}[mybox]
\begin{defn}
\label{defn-M-4.5}
\(M\) complex manifoldとし, 
\(
\varphi : M \to \Lambda^{\ell}T^{*}(M)\otimes_{\mathbb{R}}\mathbb{C}
\)をsectionとする. 

\(\varphi\)がtype $(p,q)$であるとは, 
ある開被覆$\{U_{\lambda}\}_{\lambda\in\Lambda}$で各$U_{\lambda}$は座標近傍であり. 
\(\varphi|_{U_{\lambda}}\)が$U_\lambda$上でtype$(p,q)$であることとする.  
\end{defn}
\end{tcolorbox}


\begin{rem}
\ref{rem-M-4.4}より, この定義は \(\{U_{\lambda}\}_{\lambda}\)のchoiceによらない. 
\end{rem}

\begin{tcolorbox}[mybox]
\begin{lem}
\label{lem-M-4.7}
\(M\) complex manifold, \(\dim_{\mathbb{C}} M= m\), \(0\leq k\leq 2m\)とする. 
\begin{enumerate}
\item $\C$ベクトル空間のsheafとして\(\mathcal{C}^{k}=\bigoplus_{\substack{p,q\geq 0\\ p+q=k}}\mathcal{C}^{(p,q)}\)である.
ここで,  \(\mathcal{C}^{(p,q)}\) というsheafを
\[
\mathcal{C}^{(p,q)} : U \longmapsto
 \mathcal{C}^{(p,q)}(U)=\{\text{$U$上の連続係数の$(p,q)$-form}\}
\]
として定める. 
同様に, $\C$ベクトル空間のsheafとして
\(\mathcal{E}^{k}=\bigoplus_{\substack{p,q\geq 0\\ p+q=k}}\mathcal{E}^{(p,q)} \)である. 
\item  位相$\C$ベクトル空間として\(\mathcal{K}^{k}(M)=\bigoplus_{\substack{p,q\geq 0\\ p+q=k}}\mathcal{K}^{(p,q)}(M)\), 
\(
\mathcal{D}^{k}(M)=\bigoplus_{\substack{p,q\geq 0\\ p+q=k}}\mathcal{D}^{(p,q)}(M)
\)が成り立つ. 

\end{enumerate}
\end{lem}
\end{tcolorbox}

ここで,  \(\mathcal{K}^{(p,q)}(M)\), \(\mathcal{D}^{(p,q)}(M)\) に位相を次の部分位相で入れる:
\[
\begin{aligned}
\mathcal{K}^{(p,q)}(M) &:= \mathcal{K}^{p+q}(M)\cap \mathcal{C}^{(p,q)}(M)
\quad 
\mathcal{D}^{(p,q)}(M) &:= \mathcal{D}^{p+q}(M)\cap \mathcal{C}^{(p,q)}(M)
\end{aligned}
\]

\begin{proof}

(1). 定義から, , \(\mathcal{C}^{(p,q)}\), \(\mathcal{E}^{(p,q)}\) 
は\(p+q=k\)としたときの
\(
\mathcal{C}^{k},\ \mathcal{E}^{k} 
\)
のsubsheafである. よって局所近傍をとれば直和であることが言える. 

(2). $\omega \in \mathcal{K}^{k}(M)$について, 
 \(
\omega=\sum_{p+q=k}\omega_{p, q}
\) で$\omega_{pq}\in \mathcal{C}^{(p,q)}(M)$となるようにとる.
\(
\mathrm{Supp}\,\omega=\bigcup_{p+q=k}\mathrm{Supp}\,\omega_{pq}\)
より, 
\(\omega_{pq}\in \mathcal{K}^{(p,q)}(M)\)である.
よって$\C$ベクトル空間として
\[
\mathcal{K}^{k}(M)
=\bigoplus_{p+q=k}
\mathcal{K}^{(p,q)}(M)
\]
である. ここで
\[
\mathcal{K}^{k}(M)\xrightarrow{\mathrm{pr}}
\mathcal{K}^{(p,q)}(M)\hookrightarrow \mathcal{K}^{k}(H)
\]
を考える. ($\mathrm{pr}$は連続である.)
任意の位相$\C$ベクトル空間
$X$について. 
\[
\mathrm{Hom}(\mathcal{K}^{k}(M),X)
\simeq \prod_{p+q=k}\mathrm{Hom}(\mathcal{K}^{(p,q)}(M),X)
\quad
f \longmapsto f|_{\mathcal{K}^{(p,q)}(M)}
\]
は同型になる. (単射は自明で,全射$\mathrm{pr} : \mathcal{K}^{k}(M)\to \mathcal{K}^{(p,q)}(M)$が連続なので)
よって$\mathcal{K}^{k}(M) \cong \oplus_{p+q=k} \mathcal{K}^{(p,q)}(M)$である. 
 \(\mathcal{D}^{k}(M)\)も同様である. 
\end{proof}

\begin{tcolorbox}[mybox]
\begin{defn}[{Type of currents}]
\label{defn-M-4.8}
\(M\) complex manifold, 
 \(\dim_{\mathbb{C}} M=m\)とする. 
\(T\in \mathcal{D}'^{\,\ell}(M)\) \ (\(\ell\geq 0\))について, 
\[
T:\mathcal{D}^{2m-\ell}(H)\underset{\text{as top $\C$-v.s.}}{=}
\bigoplus_{s+t=2m-\ell, s,t\geq 0} \mathcal{D}^{(s,t)}(M)
\longrightarrow \mathbb{C}\ \text{conti.}
\]
である. そこで,  \(p+q=\ell,\ p,q\geq 0\)について, 
\(T\) が type \((p,q)\) であることを
\[
T|_{\mathcal{D}^{(s,t)}(H)}=0
\quad \text{for } \forall\ (s,t)\neq (m-p,m-q) 
\text{ s.t. } s+t=2m-\ell \, s,t\geq 0
\]
として定義する. 
% [page 1]
$\mathcal{D}'^{(p,q)}(M):=\mathcal{D}'_{(m-p,m-q)}(M):=\{T\in \mathcal{D}'^{\,p+q}(M)\mid T \text{ is of type }(p,q)\}$とする. 
\end{defn}
\end{tcolorbox}

\ref{lem-M-4.7}より, 
\(
\mathcal{D}'^{\,\ell}(M)=\bigoplus_{p+q=\ell}\mathcal{D}'^{(p,q)}(M)
\)
が位相$\C$ベクトル空間としていえる. 
ここで\(\mathcal{D}'^{(p,q)}(M)\)には\(T\in \mathcal{D}'^{\,\ell}(M)\) の部分位相を入れる. 
\begin{proof}
\ref{lem-M-4.7}から$\C$ベクトル空間の同型は言える. 
また$\varphi\in \mathcal{D}^{2m-p-q}(M)$について, 
\[
ev_{\varphi_{m-p. m-q}}
:\mathcal{D}'^{\,\ell}(M)
\xrightarrow{\ \mathrm{pr}\ }\mathcal{D}'^{(p,q)}(M)
\xrightarrow{\  \mathrm{ev}_{\varphi_{m-p,m-q}}\ }
\mathbb{C}
\]
が可換であることから言える. 
\end{proof}

\begin{tcolorbox}[mybox]
\begin{lem}
\label{lem-M-4.9}
\(M\): complex manifold, \(\dim_{\mathbb{C}}H=m\), \(p,q\geq 0\)
とする. 
\[
M \underset{\text{open}}{\supset} U \longmapsto \mathcal{D}'^{(p,q)}(U)
\]
は \(\mathcal{D}'^{\,p+q}\)のsubsheafである. 
\end{lem}
\end{tcolorbox}

\begin{proof}

(1). \(T\in \mathcal{D}'^{(p,q)}(M)\), \(U\subset M\) openとする. 
\(T|_U \in \mathcal{D}'^{(p,q)}(U)\)を示す
これは\[
\xymatrix{
\mathcal{D}^{(s,t)}(H) \ar@{^{(}->}[r] & \mathcal{D}^{p+q}(H) \ar[r]^{T} &   \mathbb{C}\\
\mathcal{D}^{(s,t)}(U) \ar@{^{(}->}[r] \ar[u] & \mathcal{D}^{p+q}(U) \ar[u] \ar[ur]_{T|_{U}} & 
}
\]
が可換なのでOK.

(2).
 \(T\in \mathcal{D}'^{\,p+q}(H)\)について, ある開被覆\( \{U_{\lambda}\}_{\lambda\in\Lambda}\)で
\(T|_{U_{\lambda}}\in \mathcal{D}'^{(p,q)}(U_{\lambda}) \qquad \forall \lambda\in\Lambda,
\)
ならば
\(
T\in \mathcal{D}'^{(p,q)}(M).
\)
であることを示す. 

1の分割
\(\{\chi_{\lambda}\}_{\lambda\in\Lambda}\) をとる. 
これは\(\{\operatorname{Supp}\chi_{\lambda}\}_{\lambda\in\Lambda}\)が局所的に有限和で, $ \sum \chi_{\lambda}=1$であり, 
\(
\chi_{\lambda}\in \mathcal{D}(M),\ 0\leq \chi_{\lambda}\leq 1
\)となるものである. 
任意の \(\varphi\in \mathcal{D}^{(s,t)}(M)\)で
\(
s+t=2m-p-q, (s,t)\neq (m-p,m-q)
\)となるものについて, 

\[
T(\varphi)=
T\left(\sum_{\lambda\in\Lambda}\chi_{\lambda}\varphi\right)
=\sum_{\lambda\in\Lambda, \operatorname{supp}\chi_{\lambda}\cap \operatorname{supp}\varphi \neq \emptyset}
T|_{U_{\lambda}}(\chi_{\lambda}\varphi)
\]
である. 
よって右辺は$T|_{U_{\lambda}}\in \mathcal{D}'^{(p,q)}(U_{\lambda})$なので, $T\in \mathcal{D}'^{(p,q)}(M)$である. 
\end{proof}





\begin{tcolorbox}[mybox]
\begin{lem}
\label{lem-M-4.10}
\(U\subset \mathbb{C}^{m}\) open, \(T\in \mathcal{D}'^{\,\ell}(U)\), \(p+q=\ell\)とする. 
\(T\) がtype \((p,q)\) であることは, 
\[
T=\sum_{J\in \mathcal{J}(m;p)\atop K\in \mathcal{J}(m;q)} T_{J\bar{K}}\wedge dz^{J}\wedge d\bar{z}^{K}
\]
となる \(T_{J\bar{K}}\in \mathcal{D}'^{\,0}(U)\)が存在することと同値. 
\end{lem}
\end{tcolorbox}

\begin{proof}
\[
T=\sum_{p+q=\ell}\ \sum_{\substack{|J|=p\\|K|=q}} T_{J\bar{K}}\wedge dz^{J}\wedge d\bar{z}^{K}.
\]
とかく. 
任意の \(J'\in \{m;s\}\), \(\forall K'\{m;t\},  \varphi\in \mathcal{D}(U)\)について, 
\begin{align*}
T(\varphi\,dz^{J'}\wedge d\bar{z}^{K'})
&=
\sum_{p+q=\ell}\ \sum_{\substack{|J|=p\\|K|=q}}
T_{J\bar{K}}(\varphi\,dz^{J}\wedge d\bar{z}^{K}\wedge dz^{J'}\wedge d\bar{z}^{K'})\\
&=
T_{J^{'0}\bar{K}^{'0}}(\varphi\,dz^{J^{'0}}\wedge d\bar{z}^{K^{'0}}\wedge dz^{J'}\wedge d\bar{z}^{K'})
\end{align*}

ここで
\(J^{'0}\)を\(J^{'0}\cup J'=\{1,\dots,m\}\)となるものとする. 
以上より
\(T\) がtype \((p,q)\) であることは, 
\[
T_{J\bar{K}}=0 \ \text{ for } \forall J,\ \forall K
\text{ s.t. } (|J|,|K|)\neq (p,q)
\]
となり, 結果が従う

\end{proof}


\begin{tcolorbox}[mybox]
\begin{defn}
\label{defn-M-4.11}
\(M\) complex manifold, \(\dim_{\mathbb{C}}M=m\), \(p,q\geq 0\)とする. 

\[
\mathcal{K}'^{(p,q)}(H):=\mathcal{K}'_{(m-p,m-q)}(H):=\mathcal{K}'^{\,p+q}(H)\cap \mathcal{D}'^{(p,q)}(H)
\]
とする. ($\mathcal{K}'^{\,p+q}(H)\hookrightarrow \mathcal{D}'^{\,p+q}(H)$に注意する)
\end{defn}
\end{tcolorbox}

\begin{rem}
\label{rem-M-4.12}
\[
\mathcal{K}'^{(p,q)}(H)=\left\{T\in \mathcal{K}'^{\,p+q}(H)\ \middle|\ T|_{\mathcal{K}^{(s,t)}(H)}=0 \text{ for } 
\begin{array}{l}
(s+t)\neq (m-p,m-q) \\
s+t=2m-p-q
\end{array}
\right\}
\]
\begin{proof}
% [page 1]
\[
\mathcal{K}^{2m-p-q}(M)=\bigoplus_{s+t=2m-p-q}\mathcal{K}^{(s,t)}(M)
\underset{\text{compatible with direct sum}}{\longleftarrow} \bigoplus_{s+t=2m-p-q}\mathcal{D}^{(s+t)}(M)
\]
であるので, 
\(
\mathcal{K}'^{(p,q)}(M)\supset \mathrm{RHS} \)が言える. 

逆側の包含を示す. 
\(T\in \mathcal{K}'^{(p,q)}(M)\)とする.
$\C$ベクトル空間として,  
$$
\mathcal{K}'^{\,p+q}(M)=\bigoplus_{p'+q'=p+q} V^{(p',q')}
\quad
T \longmapsto \sum T_{p'q'}
$$
であるので, $ (p',q')\neq (p,q)$ならば
$T_{p'q'}\big|_{\mathcal{D}^{2m-p-q}(M)}=0$
である. 
よって, $(p',q')\neq (p,q)$
ならば$T_{p'q'}=0$
なので, $T \in  V^{(p,q)}$となり右辺の集合に入る. 
\end{proof}
\end{rem}



\begin{tcolorbox}[mybox]
\begin{defn}
\label{defn-M-4.13}
 \(M\): complex manifold, \(\dim_{\mathbb{C}}M=m\),  $0\leq \ell\leq 2m$, 
\(T\in \mathcal{D}'^{\,\ell}(M)\)とする. 

\[
\partial T:\mathcal{D}^{2m-\ell-1}(M)\to \mathbb{C};\qquad \varphi\mapsto (-1)^{\ell+1}T(\partial\varphi)
\]
\[
\bar{\partial}T:\mathcal{D}^{2m-\ell-1}(M)\to \mathbb{C};\qquad \varphi\mapsto (-1)^{\ell+1}T(\bar{\partial}\varphi).
\]
とする. 
 \(\partial T,\bar{\partial}T\in \mathcal{D}'^{\,\ell+1}(M)\)である.

また
\(T\in \mathcal{D}'^{\,\ell}(M)\)について, 
\(\partial \)-closed を$\partial T=0$で定める. 
\(\bar{\partial} \)-closed も$\bar{\partial} T=0$で定める. 

\end{defn}
\end{tcolorbox}

次が成り立つ. 
\begin{itemize}
\item locallyに
\(
\partial\left(\sum \varphi_{J\bar{K}}\,dz^{J}\wedge d\bar{z}^{K}\right)
=
\sum_{j}\left(\sum \frac{\partial \varphi_{J\bar{K}}}{\partial z^{j}}\,dz^{j}\right)\wedge dz^{J}\wedge d\bar{z}^{K}
\)
\item \(
\partial=\operatorname{pr}_{r+1,s}\circ d \quad \text{on } \mathcal{E}^{(r,s)}(H).
\qquad
\partial+\bar{\partial}=d
\)
\item 
\(
\bar{\partial}=\operatorname{pr}_{r,s+1}\circ d \quad \text{on } \mathcal{E}^{(r,s)}(H),
\)
\item
\(
\partial(\omega\wedge \eta)=\partial\omega\wedge \eta+(-1)^{k}\omega\wedge \partial\eta
\qquad
\omega\in \mathcal{E}^{k}(H)
\)
\item
\(
\bar{\partial}(\omega\wedge \eta)=\bar{\partial}\omega\wedge \eta+(-1)^{k}\omega\wedge \bar{\partial}\eta
\qquad
\eta\in \mathcal{E}^{\ell}(H)
\)
\end{itemize}

\(
(\partial T,\bar{\partial}T\in \mathcal{D}'^{\,\ell+1}(M)
\), であることは, $A\subset M$コンパクトについて, 
\[
\xymatrix{
\partial T: \mathcal{D}^{2m-\ell-1}(M) \ar[r]^{\partial} &
\mathcal{D}^{2m-\ell}(H) \ar[r]^{(-1)^{\ell+1}T}_{\text{conti}} &
\mathbb{C} \\
\mathcal{D}^{2m-\ell-1}_{A}(H) \ar[r]^{\partial} \ar[u] \ar[r]_{\text{conti}} &
\mathcal{D}^{2m-\ell}_{A}(H) \ar[u]
}
\]
であるから.\(\bar{\partial}T\)も同様である. 

\begin{tcolorbox}[mybox]
\begin{lem}
\label{lem-M-4.14}
\(M\) complex manifold, \(\dim_{\mathbb{C}}M=m\)とする.
\begin{enumerate}
\item \( T\in \mathcal{D}'^{(p,q)}(M)\)について, 
\( \partial T\in \mathcal{D}'^{(p+1,q)}(H), \bar{\partial}T\in \mathcal{D}'^{(p,q+1)}(M)\)
\item \(T\in \mathcal{D}'^{\,\ell}(M)\)について, 
\(dT=\partial T+\bar{\partial}T\)かつ
\[ \partial\partial T=\bar{\partial}\bar{\partial}T=\partial\bar{\partial}T+\bar{\partial}\partial T=0
\]
\end{enumerate}
\end{lem}
\end{tcolorbox}

\begin{proof}
(1). 
任意の\(\varphi\in \mathcal{D}^{(s,t)}(H)\)
で\((s,t)\neq (m-p-1,m-q), s+t=2m-p-1-q\)について
\[
(\partial T)(\varphi)
=(-1)^{p+q-1}T(\underbrace{\partial \varphi}_{\text{type }\neq (m-p,m-q) })=0.
\]
\(\bar{\partial}T\)についても同様である. 

(2). \(\mathbb{C}^{\infty}\)-form について
\[
d=\partial+\bar{\partial},\qquad \partial\partial=\bar{\partial}\bar{\partial}
=\bar{\partial}\partial+\partial\bar{\partial}=0
\]
であるため. 
\end{proof}


\begin{tcolorbox}[mybox]
\begin{lem}
\label{lem-M-4.15}
\(M\): complex manifold, \(\dim_{\mathbb{C}}M=m\)
について以下が可換になる. 

% [page 1]
\[
\xymatrix{
\mathcal{E}^{(p,q)}(H) \ar@{^{(}->}[r]^{[\cdot]} \ar[d]_{\partial} &
\mathcal{D}'^{(p,q)}(H) \ar[d]^{\partial} \\
\mathcal{E}^{(p+1,q)}(H) \ar@{^{(}->}[r]^{[\cdot]} &
\mathcal{D}'^{(p+1,q)}(H)
}
\qquad\qquad
\xymatrix{
\mathcal{E}^{(p,q)}(H) \ar@{^{(}->}[r]^{[\cdot]} \ar[d]_{\bar{\partial}} &
\mathcal{D}'^{(p,q)}(H) \ar[d]^{\bar{\partial}} \\
\mathcal{E}^{(p,q+1)}(H) \ar@{^{(}->}[r]^{[\cdot]} &
\mathcal{D}'^{(p,q+1)}(H)
}
\]
\end{lem}
\end{tcolorbox}
\begin{proof}
\(\omega\in \mathcal{E}^{(p,q)}(M)\)とする. 

 \([\omega]\in \mathcal{D}'^{(p,q)}(M)\)であることは, 任意の\(\varphi\in \mathcal{D}^{(s,t)}(M)\)で
\((s,t)\neq (m-p,m-q), s+t=2m-p-q
\)となるものについて
\[
[\omega](\varphi)
=\underbrace{\int_{M}}_{\text{複素多様体は向きづけ可能}}
\omega\wedge \varphi
\underset{\text{$2m$次でない}}{=}0
\]

次に \(\partial[\omega]=[\partial\omega]\)を示す. 
これも \(\varphi\in \mathcal{D}^{2m-p-1-q}(M)\)について
\begin{align*}
(\partial[\omega])(\varphi)
&\underset{\ref{defn-M-4.13}}{=}(-1)^{p+q-1}[\omega](\partial\varphi)\\
&\underset{\text{$[\bullet]$の定義}}{=}
= (-1)^{p+q-1}\int_{M}\omega\wedge \partial\varphi \\
&\underset{\partial(\omega\wedge \varphi)= \partial\omega\wedge \varphi + (-1)^{p+q}\omega\wedge \partial\varphi }{=}
 \underbrace{\int_{M}\partial\omega\wedge \varphi}_{=[\partial\omega](\varphi)}
 -\underbrace{\int_{M}\partial(\omega\wedge \varphi)}_{=0 by (*)}
\end{align*}

(*)についてはStokesの定理"
任意の\( \psi\in \mathcal{D}^{2m-1}(M)\)について, \(\int_{M}\partial \psi=\int_{H}\bar{\partial}\psi=0\)"である. 
これは任意の$\psi\in \mathcal{D}^{(m-1,m)}(M)$について, 
\(\ \partial\psi=(\partial+\bar{\partial})\psi=d\psi\)
なので
\[
\int_{M}\partial\psi=\int_{M}d\psi\underset{\text{普通のStokes}}{=}0
\]
である. $\bar{\partial}$も同様である. 


% [page 1]
\begin{ex}

 \( \omega \in L_{\mathrm{loc}}^{(p,q)}(M) \)とし, 
 \([ \omega ] \in \mathcal{D}'^{p+q}(M)\)について, 

\[
[\omega](\varphi):=\int_M \omega \wedge \varphi
\qquad \text{for } \varphi \in \mathcal{D}^{2n-p-q}(M)
\]
と定義すると, \([ \omega ] \in \mathcal{D}'^{(p,q)}(M)\)となる.

\begin{proof}
\(\bigl((r,s)\neq (n-p,n-q)\bigr)\)かつ
\( r+s=2n-p-q\)
とする. 任意の
\( \varphi \in \mathcal{D}^{(r,s)}(M) \)について, 
\(
\omega \wedge \varphi = 0
\)
である. よって$[\omega](\varphi)=0$
\end{proof}

\end{ex}


\begin{ex}
 \( N \subset M \)を \( \dim_{\mathbb{C}} N = k \)次元のcomplex submanifoldとする. 

この時 \([N] \in \mathcal{D}'^{2n-2k}(M)\) を
\[
[N](\varphi)=\int_N \varphi|_N
\qquad \text{for } \varphi \in \mathcal{D}^{2k}(M).
\]
と定義すると,  \([N] \in \mathcal{D}'^{(n-k,n-k)}(M)\)

\begin{proof}
\( r+s=2k \)かつ$(r,s)\neq (k,k)$とする. 
\( \varphi \in \mathcal{D}^{(r,s)}(M) \) について, 
\[
\varphi|_N = 0.
\]
である. 
これはアトラス \( \{U_\lambda; z_\lambda^1,\cdots,z_\lambda^n\}_{\lambda\in\Lambda} \)
で
\(
U_\lambda \cap N = \bigl( z_\lambda^{k+1}=\cdots=z_\lambda^n=0 \bigr).
\)
となるものをとればわかる. 
\end{proof}
\end{ex}

\subsection{複素共役}

複素共役
\( (\Lambda^k T^*M) \otimes_{\mathbb{R}} \mathbb{C}
\overset{\text{id}\otimes(-)}{\longrightarrow}
(\Lambda^k T^*M)\otimes_{\mathbb{R}} \mathbb{C} \)
によって, \( \mathcal{E}^k, \mathcal{E}^{p,q}, \mathcal{D}'(M), \mathcal{K}^k(M) \)に複素共役を誘導する. 

\begin{tcolorbox}[mybox]
\begin{defn}
\label{defn-M-4.2.1}
\( l\geq 0 \). \( T\in \mathcal{D}'^{\,l}(M) \)について, 
 \( \overline{T} \) を
\[
\overline{T}:\mathcal{D}^{n-l}(M)\to \mathbb{C};\ \varphi \mapsto \overline{T(\overline{\varphi})}.
\]
この時\( \overline{T}\in \mathcal{D}'^{\,l}(M) \)である

 \( T=\overline{T} \)となるとき, \( T \) をreal current (実カレント)という. 
\end{defn}
\end{tcolorbox}

\begin{proof}[\( \overline{T}\in \mathcal{D}'^{\,l}(M) \)の証明]
 \( \overline{T} \) が \( \mathbb{C} \)線形は明らか. 連続性は以下の図式より. 
\[
\xymatrix@C=4em@R=2em{
\overline{T}\colon
D^{m-l}(M)\ar[r]^{\overline{(\cdot)}}
&
D^{m-l}(M)\ar[r]^{T}_{\text{conti}}
&
\mathbb{C}\ar[r]^{\overline{(\cdot)}}_{\text{conti}}
&
\mathbb{C}
\\
D_A^{m-l}(M)\ar@{^{(}->}[u]\ar[r]_{\text{conti}}
&
D_A^{m-l}(M)\ar@{^{(}->}[u]
}
\]
\end{proof}


また\( T\in \mathcal{D}'^{(p,q)}(M) \)ならば \( \overline{T}\in \mathcal{D}'^{(q,p)}(M) \)である. 

\begin{proof}
\( \varphi\in \mathcal{D}^{(r,s)}(M) \), \( r+s=p+q \)について
$\varphi$の型は \( (n-p,n-q) \)でないならば,  \( \varphi \)の型は\( (n-q,n-p) \)でないので
\[
\overline{T}(\varphi)=\overline{T(\overline{\varphi})}=0
\]
\end{proof}


\begin{tcolorbox}[mybox]
\begin{lem}
\label{lem-M-4.2.2}
 \( k,l\geq 0 \)とする. 

\begin{enumerate}
\item 次は可換.
\[
\xymatrix{
L^k_{\mathrm{loc}}(M) \ar[r]^{[\cdot]} \ar[d]^{\overline{(\cdot)}}
& \mathcal{D}'^{\,k}(M) \ar[d]^{(\overline{(\cdot)}}
\\
L^k_{\mathrm{loc}}(M) \ar[r]^{[\cdot]}
& \mathcal{D}'^{\,k}(M)
}
\]

\item  \( T\in \mathcal{D}'^{\,k}(M),\ \omega\in \mathcal{E}^{\,l}(M) \)ならば
\(\overline{T\wedge \omega}=\overline{T}\wedge \overline{\omega}
\)
\end{enumerate}
\end{lem}
\end{tcolorbox}
\begin{proof}
(1). \( \omega\in L^k_{\mathrm{loc}}(M) \), \( \varphi\in \mathcal{D}^{n-k}(M) \)について

\[
\overline{[\omega]}(\varphi)
=\overline{[\omega](\overline{\varphi})}
=\overline{\int_M \omega\wedge \overline{\varphi}}
=\int_M \overline{\omega\wedge \overline{\varphi}}
=\int_M \overline{\omega}\wedge \varphi
=[\overline{\omega}](\varphi).
\]

(2) \(  \varphi\in \mathcal{D}^{2n-k-l}(M) \)について, 
\[
\overline{T\wedge \omega}(\varphi)
=\overline{T\wedge \omega (\overline{\varphi})}
=\overline{T(\omega\wedge \overline{\varphi})}
=\overline{T(\overline{\overline{\omega}\wedge \varphi})}
=\overline{T}(\overline{\omega}\wedge \varphi)
=\overline{T}\wedge \overline{\omega}(\varphi)
\]
\end{proof}

\begin{tcolorbox}[mybox]
\begin{cor}
\label{cor-M-4.2.3}
 \( U\subset \mathbb{C}^m \)開集合,  \( T\in \mathcal{D}'^{(p,q)}(U) (p,q\geq 0) \)について
$$
 T=\sum_{J \in \{m ; p\} K \in \{m ; q\} } 
 T_{J\overline{K}} \wedge dz^J \wedge d\overline{z}^K 
 $$
 となるような
 \( T_{J\overline{K}}\in \mathcal{D}'^{\,0}(U)\)をとる. 
 (ただし \((z^1,\cdots,z^m)\)を$U$の座標とする. )

この時
\[
\overline{T}
=
\sum_{J \in \{m ; p\} K \in \{m ; q\} } 
 \overline{T_{J\overline{K}}} \wedge d\overline{z}^J \wedge d z^K 
\]
\end{cor}
\end{tcolorbox}
\begin{proof}
 Lem \ref{lem-M-4.2.2} (2)より.
\end{proof}

\begin{tcolorbox}[mybox]
\begin{lem}
\label{lem-M-4.2.4}
\( k\geq 0 \)について, $T\in \mathcal{D}'^{\,k}(M)$ならば. 
\[
\overline{dT}=d\,\overline{T},
\qquad
\overline{\partial T}=\overline{\partial}\,\overline{T},
\qquad
\overline{\overline{\partial}T}=\partial \overline{T}
\]
\end{lem}
\end{tcolorbox}
\begin{proof}

\(  \varphi\in \mathcal{D}^{2n-k-1}(M) \)について, 

\[
\overline{dT}(\varphi)
=
\overline{dT(\overline{\varphi})}
=
\overline{(-1)^{k+1}T(d\overline{\varphi})}
=
(-1)^{k+1}\overline{T(\overline{d\varphi})}
=
(-1)^{k+1}\overline{T}(d\varphi)
=
(d\overline{T})(\varphi)
\]


% [page 1]
\( \partial,\overline{\partial} \) についても
\(
\overline{\partial \varphi}=\overline{\partial}\,\overline{\varphi},
\overline{\overline{\partial}\varphi}=\partial \overline{\varphi}
\)
により同様である. 
\end{proof}

\subsection{Pushforward of current}
\begin{tcolorbox}[mybox]
\begin{defn}
\label{defn-M-4.2.5}
 \( M,N \) 2nd countable complex manifoldとする. 
 \(\dim_{\mathbb{C}}M=m,\dim_{\mathbb{C}}N=n\)とする. 

正則写像\( f:M\to N \) をquasi-proper
(i.e. 任意の\(  K\subset N \) cptについて\( f^{-1}(K) \) is cptである)と仮定する. 
この時
\[
f_*:\mathcal{D}'^{\,k} (M)\to \mathcal{D}'^{\,k} (N)
\quad
T\longmapsto \left( \mathcal{D}^{k}(N)\overset{f^*}{\longrightarrow}\mathcal{D}^{k}(M)\overset{T}{\to}\mathbb{C} \right)
\]
として定義する. すると\( f_*T\in \mathcal{D}'^{\,\ell}(N) \)であり, $f_{*}$は連続である(後述)
$ f_*T$をcurrentのpushforwardという. 
\end{defn}
\end{tcolorbox}
以下\( f_*T\in \mathcal{D}'^{\,\ell}(N) \)であり, $f_{*}$は連続であることを示す. 
\begin{proof}
\[
\xymatrix@C=4em@R=2em{
D^{k}(N)\ar[r]^{f^{*}}
&
D^{k}(M)\ar[r]^{T}_{\text{conti $\C$ 線型}}
&
\mathbb{C}
\\
D_{A}^{k}(N)\ar@{^{(}->}[u]\ar[r]^{f^{*}}_{\text{conti}}
&
D_{f^{-1}(A)}^{k}(M)\ar@{^{(}->}[u]
&
}
\]
ここで$f$がquasi-properなので, $f^{-1}(A)$がコンパクトになることに注意する. 
よって\( f_*T\in \mathcal{D}'^{\,\ell}(N) \)である. 

\( f_* \) が \( \mathbb{C} \)-線形は明らか. 
\( f_* \)の連続性は, 
任意の\(\varphi\in \mathcal{D}^{\ell}(N)\)について

\[
\mathcal{D}'_{\ell}(M)\overset{f_*}{\longrightarrow}\mathcal{D}'_{\ell}(N)\overset{\mathrm{ev}_\varphi}{\longrightarrow}\mathbb{C}
\quad 
T\longmapsto f_*T \longmapsto f_{*}T[\varphi]=T[f^*\varphi]
\]
をかんがえると, 上はev\(_{f^*\varphi}\)に等しいことから. ($\mathcal{D}'_{\ell}(M)$にはevaluationが連続になるような最小の位相を入れていることに注意する.)
\end{proof}


\begin{tcolorbox}[mybox]
\begin{lem}
\label{lem-M-4.2.6}
\( f:M\to N \) quasi-proper 正則写像とする. 
この時次が可換である. 
\[
\xymatrix{
\mathcal{D}'_{k}(M) \ar[r]^{f_*} \ar[d]_{\partial}
& \mathcal{D}'_{\ell}(N) \ar[d]^{\partial}
\\
\mathcal{D}'_{\ell-1}(M) \ar[r]^{f_*}
& \mathcal{D}'_{\ell-1}(N)
}
,\qquad
\xymatrix{
\mathcal{D}'_{\ell}(M) \ar[r]^{f_*} \ar[d]_{\overline{\partial}}
& \mathcal{D}'_{\ell}(N) \ar[d]^{\overline{\partial}}
\\
\mathcal{D}'_{\ell-1}(M) \ar[r]^{f_*}
& \mathcal{D}'_{\ell-1}(N)
}
\]
特に \( d\circ f_* = f_*\circ d \).
\end{lem}
\end{tcolorbox}
 (submanifoldについては\( d\circ f_* = (-1)^{\dim_{\mathbb{R}}M+\dim_{\mathbb{R}}N} f_*\circ d \) であることに注意)

\begin{proof}
\( T\in \mathcal{D}'_{k}(M) \), \( \varphi\in \mathcal{D}^{k-1}(N) \)について, 
\begin{align*}
(d(\underbrace{f_*T}_{\in  \mathcal{D}'_{k}(N)= \mathcal{D}'^{2n-k}(N)  }))(\varphi)
&\overset{\text{def}}{=}
(-1)^{2n-k-1}(f_*T)(d\varphi)
\\
&=
(-1)^{2n-k-1}T[f^*d\varphi]
\\
&=
(-1)^{2n-k-1}\underbrace{T}_{ \in \mathcal{D}'_{k}(N)= \mathcal{D}'^{2m-k}(N)  }[df^*\varphi]
\\
&=
(-1)^{2n-k-1}(-1)^{2m-k-1}(dT)(f^*\varphi)
\\
&=
(-1)^{2k+2m}(f_*dT)(\varphi)
\end{align*}
よって
$d\circ f_* = (-1)^{2k+2m} f_*\circ d = f_*\circ d$となる. 
\( \partial,\overline{\partial} \) についても同じである. 
\end{proof}

\subsection{Distribution and measure}

\begin{tcolorbox}[mybox]
\begin{defn}
\label{defn-M-4.2.7}
$X$をlocally compact Hausdorff topological spaceとする. 
正Borel測度\( \mu \)とする. 

\( \mu \) がregularとは, 
任意のBorel 集合\( E\subset X \) について
\[
\mu(E)=\inf\{\,\mu(U)\mid E\subset U\subset X,\ U \text{ open}\,\}
=\sup\{\,\mu(K)\mid K\subset E,\ K \text{ cpt}\,\}
\]

\( \mu \) がlocally finiteとは, 
任意の
\( x\in X,\)について, ある開集合\( x\in  U\subset X \)があって, 
\(\mu(U)<\infty
\)となること. 
これは任意のコンパクト
\(K\subset X \)について\(\mu(K)<\infty\)であることと等しい. 
\end{defn}
\end{tcolorbox}


\begin{tcolorbox}[mybox]
\begin{defn}
\label{defn-M-4.2.8}
 $(X, m )$測度空間とする. 
$\mu$が複素測度であるとは, 写像
\[
\mu: m \to \mathbb{C}
\]
であって任意の
 $\forall E_i\in m\ (i=1,2,\cdots)$ で $E_i\cap E_j=\phi$ ($i\neq j$)となるものについて
 $$
 \mu\left(\bigcup_{i=1}^{\infty} E_i\right)=\sum_{i=1}^{\infty}\mu(E_i)
 $$
ただし, 右辺が常に収束することを要求する.
(特に $\sum_{i=1}^{\infty}\mu(E_i)$ は絶対収束する)

$E\in m$について. 
\[
|\mu|(E):=\sup\left\{\sum_{i=1}^{\infty}|\mu(E_i)|\ \middle|\ E=\bigcup_{i=1}^{\infty}E_i,\ 
E_i\cap E_j=\phi \ \text{for}\ i\neq j\right\}
\]
とし, 
$|\mu|$を$\mu$のtotal variationと呼ぶ

\end{defn}
\end{tcolorbox}

\begin{tcolorbox}[mybox]
\begin{thm}\cite[Theorem 6.2, 6.4]{Rud2}
\label{thm-M-4.2.9}
上の$|\mu|$は, $|\mu|(X)<\infty$となる$(X,\mathcal{M})$上のpositive measure である. 
\end{thm}
\end{tcolorbox}

\begin{tcolorbox}[mybox]
\begin{defn} 
\label{defn-M-4.2.10}
$X$をlocally compact Hausdorff topological spaceとする. 
複素Borel測度$\mu$について, $|\mu|$ がregularである時, 
$\mu$ が regularであると定義する. 
\end{defn}
\end{tcolorbox}

以下\cite{Rud2}で使うFactを並べる. 
\begin{tcolorbox}[mybox]
\begin{thm}\cite[Theorem 2.18]{Rud2}
\label{thm-M-4.2.11}
$X$をlocally compact Hausdorff topological spaceとする. 
さらに任意の開集合が$\sigma$-コンパクトであると仮定する.
この時任意のlocally finite positive Borel 測度はregular.
\end{thm}
\end{tcolorbox}
$\sigma$-コンパクトとは可算個のコンパクト集合の合併になること. 2nd conutableなら$\sigma$コンパクト

次はリースの表現定理である. 
\begin{tcolorbox}[mybox]
\begin{thm}\cite[Theorem 2.14]{Rud2}
\label{thm-M-4.2.12}
$X$をlocally compact Hausdorff topological spaceとする. 
さらに任意の開集合が$\sigma$-コンパクトであると仮定する.

$T:\mathcal{K}(X)\to \mathbb{C}$を $\mathbb{C}$線型写像とし, さらにpositiveであると仮定する. 
 ("$T$がpositive"とは任意の$ \varphi\in \mathcal{K}(X)$で$\varphi \ge 0$なるものについて, $(T(\varphi)\geq 0 $であること)
 
この時positive locally finite regular Borel 測度 $\mu$であって
\[
T(\varphi)=\int_X \varphi \, d\mu \quad \text{for } {}^\forall \varphi\in \mathcal{K}(X).
\]
となるものがただ一つ存在する, 
\end{thm}
\end{tcolorbox}
ここで
\(
\left( \mathcal{K}(X)=\{\varphi:X\to \mathbb{C}\ \mathrm{cont.};\ \mathrm{supp}\ \varphi\ \mathrm{cpt}\}\right)
\)である. \cite{NO}の記号に合わせた. 

% [page 1]


\begin{tcolorbox}[mybox]
\begin{thm}\cite[Theorem 6.19]{Rud2}
\label{thm-M-4.2.13}
$X$をlocally compact Hausdorff topological spaceとする. 
$\C$線型写像$T:\mathcal{K}(X)\to \mathbb{C}$で

ある$ C>0$があって, 
$$\ |T(\varphi)|\leq C \sup_{x\in X}|\varphi(x)| \quad \text{for } {}^\forall \varphi\in \mathcal{K}(X),
$$
ならば
complex regular Borel 測度 $\mu$であって
\[
T(\varphi)=\int_X \varphi \, d\mu \qquad \text{for } {}^\forall \varphi\in \mathcal{K}(X)
\]
となるものがただ一つ存在する. 

さらにこの場合
\[
\|T\|:=\sup\left\{\frac{|T(\varphi)|}{\sup_{x\in X}|\varphi(x)|}\ \middle|\ \varphi\in \mathcal{K}(X),\ \varphi\neq 0\right\}
=|\mu|(X).
\]

\end{thm}
\end{tcolorbox}

\begin{rem}
$X$を$C^{\infty}$ 2nd countable manifoldとする. 
%$T$がpositiveという条件は強く, これから$T$が連続であることが従う. 
$T:\mathcal{K}(X)\to \mathbb{C}$を$\mathbb{C}$線形かつpositive とする.
ここでpositiveとは
\[
\varphi\geq 0 \Rightarrow T(\varphi)\geq 0
\]
を意味する. すると$T$は連続になる. もっと強くorder $0$ distributionである. 
\begin{proof}
$A\subset X$をコンパクト集合とする. 
 $\psi\in \mathcal{D}(X)$で, $\psi\geq 0$かつ$A$上で$\psi \equiv 1$となるものを取る. 
任意の$\varphi\in \mathcal{D}_A(X)$について, 
もし$\varphi(X)\subset \mathbb{R}$ならば
\[
\underbrace{\|\varphi\|_0}_{=:\sup_{x\in X} |\varphi(x)|}\psi \pm \varphi \geq 0
\]
であるので, $T$がpositiveから
\[
T(\|\varphi\|_0\psi \pm \varphi)\geq 0
\]
以上より線形性から, 
\( |T(\varphi)|\leq T(\psi)\|\varphi\|_0
\)を得る.

一般の場合($\varphi(X)\subset \mathbb{C}$の場合), 

\[
|T(\varphi)|=\left|T(\Re\varphi)+iT(\Im\varphi)\right|
\leq T(\psi)(\|\Re\varphi\|_0+\|\Im\varphi\|_0)
\underset{\text{上の議論}}{\leq} 2T(\psi)\|\varphi\|_0.
\]
よって$C_A:=2T(\psi)$ととれば, $T|_{\mathcal{D}(X)}$はorder $0$ distributionの定義を満たす. 
\end{proof}
\end{rem}

\end{proof}



